\documentclass[11pt,a4paper]{article}
\usepackage[utf8]{inputenc}
\usepackage[T1]{fontenc}
\usepackage[italian]{babel}
\usepackage{amsmath,amssymb,amsthm,mathtools}
\usepackage[margin=26mm]{geometry}
\usepackage[colorlinks=true,linkcolor=blue,urlcolor=blue,citecolor=blue]{hyperref}
\newtheorem{theorem}{Teorema}
\newtheorem{lemma}[theorem]{Lemma}
\newtheorem{corollary}[theorem]{Corollario}
\newtheorem{proposition}[theorem]{Proposizione}
\theoremstyle{remark}
\newtheorem*{remark}{Osservazione}
\newcommand{\R}{\mathbb R}
\newcommand{\Sph}{\mathbb S^2}
\newcommand{\ind}{\mathbf 1}
\newcommand{\vol}{\operatorname{Vol}}
\newcommand{\supp}{\operatorname{supp}}
\newcommand{\tr}{\operatorname{tr}}
\title{Covariogramma in $\R^3$:\\proposta di unicità con regolarità del primo corpo}
\author{Quaderno di lavoro --- deduzioni e verifiche della sessione}
\date{9 ottobre 2026}
\begin{document}
\maketitle

\begin{abstract}
Si presenta una proposta di dimostrazione dell'unicità del covariogramma
per un corpo $K\in C^9_+$ in $\R^3$, confrontato con un corpo convesso
$L$ arbitrario, modulo traslazioni e riflessioni centrali. È il caso
segnalato come aperto da Bianchi dopo il suo teorema per due corpi
regolari. La proposta combina due passaggi: ricavare quasi ovunque
le coppie di densità di area del concorrente mediante il supporto
$C^{1,1}$ e gli integrali di Laplace; ordinare globalmente tali coppie
mediante i rami analitici di zeri descritti da Bianchi e la
fattorizzazione di Hadamard. Per il caso $C^9_+$ la dimostrazione è
indipendente dagli sviluppi accessori che abbassano la regolarità.
Il quaderno conserva inoltre criteri di bilanciamento, esempi di scambi
locali e una condizione di Dini sulle sezioni. La proposta principale
è accompagnata dalle dimostrazioni dei due passaggi e da controlli
dei limiti; non è presentata come un risultato già pubblicato o
convalidato esternamente. Il problema per un primo corpo convesso
arbitrario non viene risolto.
\end{abstract}

\section{Ambito e fatti esterni utilizzati}
\paragraph{Obiettivo aggiornato su indicazione dell'utente.}
È ora sufficiente risolvere un caso ancora aperto imponendo regolarità
alla frontiera. Il primo obiettivo preciso è quindi il seguente:
\begin{equation}\label{eq:target}
K\in C^9_+,\quad L\text{ corpo convesso in }\R^3,\quad g_K=g_L
\quad\Longrightarrow\quad L=v+K\text{ oppure }L=v-K.
\end{equation}
Qui non si impone alcuna regolarità a $L$. Questo distingue il problema
dal risultato già noto per due corpi $C^9_+$: \cite[Teorema 3.25 e
paragrafo immediatamente successivo, p.~23 del manoscritto]{B23}
indica esplicitamente come aperta la determinazione fra tutti i corpi
convessi. Le ricerche bibliografiche aggiornate al 9 ottobre 2026 non
hanno individuato una risoluzione successiva di questo passaggio;
questa ricerca non è una certificazione esaustiva della novità.
La proposta della Sezione \ref{sec:transport} mira a
\eqref{eq:target}, e anzi a una regolarità inferiore. La sua verifica
matematica rimane distinta dall'accertamento bibliografico.
Il 9 ottobre l'utente ha confermato che la dimostrazione di questo
preciso caso aperto costituisce il traguardo conclusivo della ricerca.
Le estensioni a regolarità inferiore presenti nel quaderno sono
sviluppi accessori e non sono necessarie a tale traguardo.

\paragraph{Versione consolidata per il riesame.}
Su richiesta dell'utente, questa versione registra la catena completa
della dimostrazione candidata di \eqref{eq:target}, con i passaggi
geometrici e analitici sviluppati nel testo e i risultati esterni
identificati. Nei controlli interni svolti non è stata individuata
una lacuna nella catena principale. Questo giudizio non equivale a
una verifica indipendente né a una certificazione della novità.
La registrazione del documento è stata chiusa per consentire
all'utente di rifare il punto. Su successiva richiesta è stato
avviato il controllo formale parziale descritto di seguito.

\paragraph{Stato del controllo Lean, 9 ottobre 2026.}
Sono stati formalizzati e controllati dal kernel 33 lemmi di supporto
e deduzioni condizionali, senza lacune ammesse o assiomi matematici
aggiunti. Il teorema \eqref{eq:target} non è ancora formalizzato né
certificato. In particolare, le identità geometriche e spettrali
assunte in alcuni lemmi Lean non sono ancora ricavate formalmente
dall'uguaglianza dei covariogrammi. La Sezione \ref{sec:leancheck}
specifica il perimetro del controllo; il codice e i registri sono
nella cartella \texttt{formalizzazione} del progetto.

\paragraph{Percorso breve per l'obiettivo attuale.}
Il lettore interessato a \eqref{eq:target} può seguire il Lemma
\ref{lem:roughcompetitor}, che recupera le coppie di densità di area
per un concorrente arbitrario; il Lemma \ref{lem:mean}, che ne rende
globale l'ordinamento; e il Corollario \ref{cor:targetC9}, che combina
questi passaggi con gli zeri descritti in \cite{B16} e l'unicità di
Minkowski. Il criterio di Dini serve all'estensione a regolarità
inferiore. Il testo contiene una proposta di dimostrazione di questa
catena di implicazioni; non è presentato come un teorema già
pubblicato o sottoposto a una revisione esterna.

Un corpo convesso è compatto e ha interno non vuoto. Il suo covariogramma è
\[
g_K(x)=\int_{\R^3}\ind_K(y)\ind_K(y-x)\,dy.
\]
Traslazioni e riflessioni centrali non lo modificano. La classe $C^2_+$
consiste di corpi con bordo $C^2$ e curvatura gaussiana positiva.
Per $u\in\Sph$, poniamo
\[
f_K(u)=\frac1{\kappa_K(u)},\qquad dS_K(u)=f_K(u)\,d\sigma(u),
\]
dove la curvatura è valutata nel punto con normale esterna $u$.

I fatti esterni centrali sono i seguenti.
\begin{enumerate}
\item Per corpi $C^2_+$, il covariogramma determina la coppia non ordinata
      $\{f_K(u),f_K(-u)\}$ per ogni $u$; si veda \cite[Proposizione 3.1]{B16}.
\item La misura di area determina un corpo convesso a meno di traslazione
      (unicità di Minkowski). Inoltre
      \begin{equation}\label{eq:balance}
      \int_{\Sph}u f_K(u)\,d\sigma(u)=0.
      \end{equation}
      Quest'ultima identità segue direttamente dal teorema della divergenza
      applicato a ogni campo vettoriale costante.
\end{enumerate}
Per il primo fatto basta conoscere il covariogramma vicino all'origine e ai
punti di contatto sul bordo di $K-K$. Il criterio della prossima sezione usa
solo la coppia non ordinata e \eqref{eq:balance}, non ulteriori valori interni.

\section{Un criterio di unicità mediante bilanciamento}
Definiamo la funzione dispari continua
\[
d(u)=f_K(u)-f_K(-u),\qquad U_+=\{u:d(u)>0\}.
\]
Supponiamo inizialmente che $U_+$ abbia un numero finito $m\geq1$ di
componenti connesse $C_1,\ldots,C_m$, e definiamo
\begin{equation}\label{eq:vi}
v_i=\int_{C_i}u\,d(u)\,d\sigma(u)\in\R^3.
\end{equation}
Le componenti di $U_-=\{d<0\}$ sono $-C_i$.

\begin{lemma}\label{lem:selection}
Se $K,L\in C^2_+$ e $g_K=g_L$, esiste un sottoinsieme
$I\subseteq\{1,\ldots,m\}$ tale che
\[
f_L(u)=
\begin{cases}
f_K(-u),&u\in\displaystyle\bigcup_{i\in I}(C_i\cup(-C_i)),\\
f_K(u),&\text{altrove}.
\end{cases}
\]
Necessariamente
\begin{equation}\label{eq:zerosum}
\sum_{i\in I}v_i=0,\qquad \sum_{i=1}^m v_i=0.
\end{equation}
\end{lemma}
\begin{proof}
Su ciascun $C_i$ i due valori $f_K(u),f_K(-u)$ sono distinti. La funzione
\[
\frac{f_L(u)-f_K(-u)}{f_K(u)-f_K(-u)}
\]
è continua e assume solo i valori $0,1$; è quindi costante. La coppia non
ordinata forza la scelta corrispondente su $-C_i$. Su $\{d=0\}$ non vi è
ambiguità. Posto $V=\bigcup_{i\in I}(C_i\cup(-C_i))$, abbiamo
$f_L-f_K=-d\ind_V$. Per la disparità di $d$,
\[
\int_{-C_i}u d(u)\,d\sigma(u)=\int_{C_i}u d(u)\,d\sigma(u)=v_i.
\]
Sottraendo le due identità \eqref{eq:balance} otteniamo
$0=-2\sum_{i\in I}v_i$. Infine
\[
2\sum_i v_i=\int_{\Sph}u d(u)\,d\sigma(u)
=2\int_{\Sph}u f_K(u)\,d\sigma(u)=0.
\]
\end{proof}

\begin{theorem}[Criterio di bilanciamento]\label{thm:balance}
Sia $K\in C^2_+$, con le componenti e i vettori di \eqref{eq:vi}.
Se nessun sottoinsieme proprio non vuoto degli indici ha somma vettoriale
nulla, allora per ogni $L\in C^2_+$
\[
g_L=g_K\quad\Longrightarrow\quad L=a+K\ \text{oppure}\ L=a-K.
\]
\end{theorem}
\begin{proof}
Il lemma impone $I=\varnothing$ oppure $I=\{1,\ldots,m\}$.
Pertanto $S_L=S_K$ oppure $S_L=S_{-K}$. L'unicità di Minkowski conclude.
\end{proof}

\begin{corollary}\label{cor:smallm}
Il criterio vale nei seguenti casi:
\begin{enumerate}
\item $m=1$;
\item $m=2$ e $v_1\ne0$ (allora $v_2=-v_1$);
\item $m=3$ e tutti i $v_i$ sono non nulli.
\end{enumerate}
Nel caso $m=2$, è sufficiente che una componente sia contenuta in un
emisfero aperto. Nel caso $m=3$, è sufficiente che ciascuna componente
sia contenuta in un emisfero aperto, eventualmente diverso.
\end{corollary}
\begin{proof}
Per $m=3$, una somma nulla di due vettori equivale all'annullarsi del terzo.
Se $C_i\subset\{u:a\cdot u>0\}$, allora
$a\cdot v_i=\int_{C_i}(a\cdot u)d(u)\,d\sigma(u)>0$.
\end{proof}

Se $d\equiv0$, non ci sono componenti e la coppia non ordinata impone
direttamente $f_L=f_K$. Anche questo caso dà unicità.
Se esistono $b$ sottoinsiemi bilanciati, compresi quello vuoto e quello
totale, il numero dei possibili concorrenti $C^2_+$, modulo traslazioni e
riflessioni, è al più $b/2$: ciascuna scelta determina al più una misura
di area e quindi al più un corpo modulo traslazione. Non si afferma che
ogni scelta bilanciata dia lo stesso covariogramma, lo stesso corpo
differenza, o un corpo $C^2_+$.

\subsection*{Una famiglia esplicita a due componenti}
Sia $z=u_3$, $P_3(z)=(5z^3-3z)/2$ e
\[
h_K(u)=1-\eta P_3(u_3),\qquad 0<\eta<1/5.
\]
Per una funzione $h$ sulla sfera poniamo $Q(h)=\nabla^2_{\Sph}h+hI$.
Gli autovalori di $Q(P_3)$, nelle direzioni meridiana e parallela, sono
\[
15z-20z^3,\qquad -5z^3.
\]
Entrambi sono compresi tra $-5$ e $5$. Quindi $Q(h_K)>0$ e $h_K$ è
la funzione di supporto di un corpo $C^\infty_+$. In dimensione tre,
$f_K=\det Q(h_K)$; usando $\tr Q(P_3)=-10P_3$ otteniamo esattamente
\[
d(u)=20\eta P_3(u_3).
\]
Con $a=\sqrt{3/5}$, le due componenti positive sono
$C_1=\{z>a\}$ e $C_2=\{-a<z<0\}$. Per simmetria assiale,
\[
v_1=40\pi\eta\left(\int_a^1 zP_3(z)\,dz\right)e_3
=8\pi\eta(3/5)^{3/2}e_3\ne0.
\]
Il Teorema \ref{thm:balance} determina quindi questo $K$ fra tutti i
concorrenti $C^2_+$, senza imporre loro regolarità $C^9$.
L'esempio verifica che il criterio non è vuoto; non ne stabilisce la novità.

\section{Una rigidità esatta della parte pari}
Per una funzione reale limitata e a supporto compatto $f$ definiamo
\[
\widetilde f(x)=f(-x),\quad E_f=(f+\widetilde f)/2,\quad
O_f=(f-\widetilde f)/2.
\]
La commutatività della convoluzione dà l'identità
\begin{equation}\label{eq:eo}
f*\widetilde f=E_f*E_f-O_f*O_f.
\end{equation}

\begin{lemma}[Assenza di divisori dello zero]\label{lem:zero}
Se $a,b\in L^1(\R^n)$ hanno supporto compatto e $a*b=0$, allora
$a=0$ oppure $b=0$, quasi ovunque.
\end{lemma}
\begin{proof}
Le trasformate di Fourier sono reali analitiche su $\R^n$ e
$\widehat a\widehat b=0$. Se $a\ne0$, l'iniettività di Fourier e la
continuità forniscono un aperto non vuoto dove $\widehat a\ne0$.
Lì $\widehat b=0$; l'unicità del prolungamento reale analitico dà
$\widehat b=0$ su tutto $\R^n$. Quindi $b=0$.
\end{proof}

\begin{theorem}\label{thm:even}
Siano $K,L\subset\R^n$ misurabili e limitati. Se
\[
g_K=g_L,\qquad
\ind_K(x)+\ind_K(-x)=\ind_L(x)+\ind_L(-x)\quad\text{q.o.},
\]
allora $K=L$ oppure $K=-L$ a meno di insiemi nulli.
Per corpi convessi l'uguaglianza è esatta.
\end{theorem}
\begin{proof}
Le parti pari sono uguali. Da \eqref{eq:eo},
\[
0=O_K*O_K-O_L*O_L=(O_K-O_L)*(O_K+O_L).
\]
Il Lemma \ref{lem:zero} impone $O_K=O_L$ oppure $O_K=-O_L$.
Nel primo caso le indicatrici coincidono; nel secondo sono riflesse.
Due corpi convessi con uguali indicatrici quasi ovunque coincidono, perché
ciascuno è la chiusura del proprio interno.
\end{proof}

Il centro può essere qualsiasi punto fissato, dopo una traslazione.
Per corpi contenenti l'origine, l'ipotesi sulle parti pari equivale a
conoscere le coppie radiali non ordinate
$\{\rho_K(u),\rho_K(-u)\}$, oppure sia $K\cup(-K)$ sia $K\cap(-K)$.
\emph{Non} equivale a conoscere $K-K$, che determina invece la parte pari
della funzione di supporto.

Se due insiemi differiscono per la riflessione di un gruppo di porzioni
attorno al medesimo centro, la loro parte pari può rimanere identica.
In tal caso il teorema esclude ogni scambio non globale con covariogramma
invariato. In particolare, se le indicatrici sono
\[
\ind_K=E+A+B,\qquad \ind_L=E+A-B,
\]
con $E$ pari, $A,B$ dispari e non nulli, allora
\begin{equation}\label{eq:cross}
g_K-g_L=-4A*B\ne0.
\end{equation}
Questa identità individua esattamente il termine di interazione che
rileva lo scambio parziale. La convessità non è necessaria al teorema.

\section{Due corpi convessi espliciti e un punto che li distingue}
Costruiamo una famiglia con regolarità finita e verifichiamo tutte le
disuguaglianze. Poniamo
\[
c=\frac9{10},\qquad d_0=\frac{19}{20},\qquad 0<\varepsilon\le10^{-5},
\]
e definiamo la funzione $C^3$
\[
\psi(t)=
\begin{cases}
0,&t\le0,\\
35t^4-84t^5+70t^6-20t^7,&0<t<1,\\
1,&t\ge1.
\end{cases}
\]
Essa è monotona, poiché $\psi'(t)=140t^3(1-t)^3$ su $(0,1)$.
Siano
\[
\beta(s)=\psi\left(\frac{s-c}{d_0-c}\right),\qquad
\phi_i(u)=\beta(u_i)-\beta(-u_i),\quad i=1,2.
\]
Le quattro calotte $\{\pm u_i>c\}$ sono disgiunte perché $2c^2>1$.
Definiamo le funzioni radiali
\[
r_+(u)=1+\varepsilon(\phi_1(u)+\phi_2(u)),\qquad
r_-(u)=1+\varepsilon(\phi_1(u)-\phi_2(u)),
\]
e i corpi stellati $K_\pm=\{tu:u\in\Sph,\ 0\le t\le r_\pm(u)\}$.

\subsection*{Verifica della convessità}
Se $r>0$ è $C^2$, un criterio sufficiente per la convessità stretta del
corpo radiale, con curvatura positiva, è la positività di
\begin{equation}\label{eq:radial}
M_r=r^2 I+2\nabla r\otimes\nabla r-r\nabla^2r.
\end{equation}
Infatti il gauge omogeneo $F(x)=|x|/r(x/|x|)$ ha Hessiana tangenziale
proporzionale a
$Q(1/r)=M_r/r^3$. La positività rende $F$ convesso e $\{F\le1\}$ convesso;
la stessa matrice dà la seconda forma fondamentale del bordo, a meno di
un fattore positivo.

Su $(0,1)$ valgono $|\psi'|\le35/16$ e $|\psi''|\le105/4$.
Dato che $|\nabla u_i|\le1$ e $\nabla^2u_i=-u_i I$, per
$\Phi=\phi_1\pm\phi_2$ otteniamo
\[
|\Phi|\le1,\quad |\nabla\Phi|\le G=\frac{175}{4},\quad
\|\nabla^2\Phi\|\le H=\frac{42175}{4}.
\]
Si usa il fatto che in ogni punto è attiva al più una calotta, e le
derivate fino all'ordine tre si raccordano. Quindi
\[
M_{r_\pm}\succeq
\left((1-\varepsilon)^2-(1+\varepsilon)\varepsilon H\right)I
\succeq\frac{35781657829}{40000000000}I>0.
\]
I due corpi appartengono a $C^3_+$; per questa scelta di $\psi$ non si
assume regolarità $C^9$. Entrambi contengono $B_{1-\varepsilon}$.

\subsection*{La differenza dei covariogrammi}
Definiamo i gusci aperti
\[
T_i=\{tu:u_i>c,\ 1-\varepsilon\beta(u_i)<t<1+\varepsilon\beta(u_i)\},
\quad a_i=\ind_{T_i},\quad A_i=(a_i-\widetilde a_i)/2.
\]
Le indicatrici di $K_+,K_-$ hanno la stessa parte pari $E$ e, quasi ovunque,
\[
\ind_{K_+}=E+A_1+A_2,\qquad
\ind_{K_-}=E+A_1-A_2.
\]
Pertanto
\begin{equation}\label{eq:explicitdiff}
g_{K_+}-g_{K_-}
=-a_1*a_2-\widetilde a_1*\widetilde a_2
 +a_1*\widetilde a_2+\widetilde a_1*a_2.
\end{equation}

\begin{proposition}[Una traslazione esplicita]\label{prop:witness}
Per $x_*=e_1+e_2=(1,1,0)$ vale
\[
g_{K_+}(x_*)-g_{K_-}(x_*)\le-\frac\pi6\varepsilon^3<0.
\]
\end{proposition}
\begin{proof}
Scriviamo $s=\sqrt{1-c^2}$. Un punto di $T_1$ ha prima coordinata
maggiore di $(1-\varepsilon)c$ e seconda coordinata di modulo al più
$(1+\varepsilon)s$; per $T_2$ scambiamo gli indici.
Poiché
\[
(1+\varepsilon)s-(1-\varepsilon)c<0,
\]
$x_*$ non appartiene a $T_1-T_2$, $-T_1+T_2$ o $-T_1-T_2$.
Gli ultimi tre termini pertinenti in \eqref{eq:explicitdiff} sono quindi
nulli in $x_*$, e resta
\[
g_{K_+}(x_*)-g_{K_-}(x_*)=-\vol\bigl(T_1\cap(x_*-T_2)\bigr).
\]
Poniamo $\delta=\varepsilon/2$. Se $y\in B_\delta(e_1)$, allora
$1-\delta<|y|<1+\delta$ e
\[
\frac{y_1}{|y|}\ge\frac{1-\delta}{1+\delta}>d_0.
\]
Quindi $\beta(y_1/|y|)=1$ e $y\in T_1$. Analogamente
$x_*-y\in B_\delta(e_2)\subset T_2$.
L'intersezione contiene $B_\delta(e_1)$, di volume
$4\pi\delta^3/3=\pi\varepsilon^3/6$.
\end{proof}

\subsection*{Perché i dati vicini agli estremi non vedono lo scambio}
Sia $\theta=\arccos c<\pi/4$. Le direzioni di due vettori appartenenti
rispettivamente a $\pm T_1,\pm T_2$ formano un angolo nell'intervallo
$[\pi/2-2\theta,\pi/2+2\theta]$. Le loro lunghezze differiscono da uno
al più di $\varepsilon$. Ne segue che ogni somma appartiene alla corona
\[
\rho_\varepsilon\le |x|\le R_\varepsilon,
\quad
\rho_\varepsilon=2\sin(\pi/4-\theta)-2\varepsilon,
\quad
R_\varepsilon=2\cos(\pi/4-\theta)+2\varepsilon.
\]
Per \eqref{eq:explicitdiff},
\begin{equation}\label{eq:annulus}
g_{K_+}(x)=g_{K_-}(x)
\quad\text{se }|x|<\rho_\varepsilon\ \text{oppure }|x|>R_\varepsilon.
\end{equation}
Uniformemente per $0<\varepsilon\le10^{-5}$,
\[
\rho_\varepsilon>0.6563,\qquad R_\varepsilon<1.8893
<1.99998\le2(1-\varepsilon).
\]
Questi decimali illustrano disuguaglianze certificate anche con stime
razionali nel file di verifica allegato.

Poiché $K_\pm$ contengono $B_{1-\varepsilon}$, i due covariogrammi sono
positivi su $B_{2(1-\varepsilon)}$. Fuori da questa palla coincidono per
\eqref{eq:annulus}, dunque hanno lo stesso supporto
$D=K_+-K_+=K_--K_-$. Il bordo di $D$ è contenuto in
$\{|x|\ge2(1-\varepsilon)\}$, dove i due covariogrammi coincidono in un
intorno. Coincidono anche vicino all'origine, ma sono distinti in $x_*$.
La stretta disuguaglianza esclude che i corpi siano traslati o riflessi
centralmente l'uno dell'altro.

\begin{remark}
L'esistenza di corpi $C^2_+$ non equivalenti con covariogrammi uguali
vicino a $\{0\}\cup\partial D$ è già affermata in
\cite[Teorema 6, ultima parte]{AB07}. La costruzione qui serve a rendere
espliciti una famiglia a regolarità finita, il termine che separa i dati,
il punto $x_*$ e una stima stretta. Non è un controesempio all'unicità
dal covariogramma completo.
\end{remark}

\section{Una ridondanza specifica della dimensione tre}
Il volume non fornisce un vincolo indipendente una volta noti sia il
corpo differenza sia le aree di tutte le proiezioni. Più precisamente,
per ogni corpo convesso in $\R^3$, con $D=K-K$,
\begin{equation}\label{eq:volume}
\vol(K)=\frac14\left(
\int_{\Sph}h_D(u)\,d(S_K+S_{-K})(u)-\vol(D)\right).
\end{equation}
Per verificarlo, scriviamo $V=\vol(K)$ e $M=V(K,K,-K)$, usando i volumi
misti. La riflessione e la simmetria degli argomenti danno
$V(K,-K,-K)=M$, quindi $\vol(D)=2V+6M$. Inoltre, essendo $h_D$ pari,
\[
\int h_D\,d(S_K+S_{-K})=2\int h_D\,dS_K
=6V(K,K,D)=6V+6M.
\]
La sottrazione dà \eqref{eq:volume}.
Le aree delle proiezioni determinano $S_K+S_{-K}$ mediante
l'iniettività della trasformata coseno sulle misure pari. Questa
osservazione non produce unicità, ma evita di contare il volume come
un'informazione ulteriore dopo questi dati.

\section{Rami analitici di zeri sotto una condizione sulle sezioni}
\label{sec:zeros}
L'idea di utilizzare l'analiticità dei rami di zeri per rendere globale
la scelta delle curvature è di \cite{B16}. Qui ne isoliamo una condizione
sufficiente in dimensione tre: occorrono soltanto zeri in una striscia
orizzontale fissata. Non si richiede una classificazione di tutti gli zeri.
La dimostrazione seguente e il successivo criterio di Dini sono sviluppi
del quaderno, ancora privi di revisione indipendente.

Per $K\in C^2_+$ poniamo
\[
\alpha_K(u)=-h_K(-u),\quad \beta_K(u)=h_K(u),\quad
w_K=\beta_K-\alpha_K,
\qquad A_K(u,t)=\mathcal H^2(K\cap\{x:x\cdot u=t\}).
\]
L'area è nulla agli estremi. Le sue derivate unilaterali sono
\begin{equation}\label{eq:slopes}
a_K(u)=\partial_t A_K(u,\alpha_K(u)+)
=\frac{2\pi}{\sqrt{\kappa_K(-u)}},\qquad
-\partial_t A_K(u,\beta_K(u)-)=b_K(u)
=\frac{2\pi}{\sqrt{\kappa_K(u)}}.
\end{equation}
Queste formule seguono anche dal calcolo locale della prossima sezione.

\paragraph{Condizioni sulle sezioni.}
Diremo che vale (S) se, per ogni $u$, la funzione
$\partial_t A_K(u,\cdot)$ è assolutamente continua sul proprio intervallo
chiuso e la sua derivata seconda interna è integrabile. Definiamo
\[
T_{K,u}(t)=\ind_{(\alpha_K(u),\beta_K(u))}(t)\,
\partial_t^2 A_K(u,t).
\]
Gli atomi dovuti alle derivate agli estremi \emph{non} fanno parte di
$T_{K,u}$. Diremo che vale (U) se vale (S) e la famiglia
$\{T_{K,u}:u\in\Sph\}$ è relativamente compatta in $L^1(\R)$.

\begin{lemma}[Asintotica in una striscia]\label{lem:strip}
Sia $K\in C^2_+$ e valga (U). Per
$\mathcal F_K(z)=\int_K e^{iz\cdot x}\,dx$ si ha
\begin{equation}\label{eq:strip}
\zeta^2\mathcal F_K(\zeta u)=
-a_K(u)e^{i\alpha_K(u)\zeta}-b_K(u)e^{i\beta_K(u)\zeta}+o(1)
\end{equation}
quando $\operatorname{Re}\zeta\to+\infty$, uniformemente in $u\in\Sph$
e in ogni striscia $|\operatorname{Im}\zeta|\leq M$ fissata.
Sotto la sola condizione (S), la stessa conclusione vale per ogni $u$
fissato, uniformemente nella parte immaginaria limitata.
\end{lemma}
\begin{proof}
Fubini e due integrazioni per parti, giustificate da (S), danno l'identità
esatta
\begin{equation}\label{eq:twoparts}
\zeta^2\mathcal F_K(\zeta u)=
-a_K(u)e^{i\alpha_K(u)\zeta}-b_K(u)e^{i\beta_K(u)\zeta}
-\int_\R T_{K,u}(t)e^{it\zeta}\,dt.
\end{equation}
Tutti i supporti sono contenuti in un intervallo comune limitato.
La chiusura della famiglia
$T_{K,u}(t)e^{-\sigma t}$, $u\in\Sph$, $|\sigma|\leq M$, è compatta
in $L^1$: è l'immagine continua di un prodotto di due compatti.
Il lemma di Riemann--Lebesgue è uniforme su un compatto di $L^1$,
come si verifica coprendolo con una rete finita di palle $L^1$ e usando
$|\widehat f|\leq\|f\|_1$. Il resto in \eqref{eq:twoparts} tende dunque
uniformemente a zero. Per $u$ fissato basta il compatto ottenuto variando
soltanto $\sigma$.
\end{proof}

\begin{lemma}[Un ramo globale per ogni indice grande]\label{lem:branches}
Sotto le ipotesi uniformi del Lemma \ref{lem:strip}, esiste $m_0$ tale che
per ogni intero $m\geq m_0$ esiste una funzione reale analitica
$z_{m,K}:\Sph\to\mathbb C$ con
$\mathcal F_K(z_{m,K}(u)u)=0$ e
\begin{equation}\label{eq:rootasymptotic}
z_{m,K}(u)=\frac{(2m+1)\pi}{w_K(u)}+i\gamma_K(u)+o(1),\qquad
\gamma_K(u)=\frac{\log(b_K(u)/a_K(u))}{w_K(u)}.
\end{equation}
Il resto è uniforme in $u$, e ciascuno zero è semplice come zero nella
variabile scalare $\zeta$.
\end{lemma}
\begin{proof}
Scriviamo $a,b,\alpha,\beta,w$ senza pedice. Gli zeri del modello
$H_u(\zeta)=a e^{i\alpha\zeta}+b e^{i\beta\zeta}$ sono
\[
\lambda_m(u)=\frac{(2m+1)\pi}{w(u)}+i\gamma(u),\qquad m\in\mathbb Z.
\]
Fissiamo $0<\delta<\pi/(4\max w)$. Sulla circonferenza
$\zeta=\lambda_m(u)+\xi$, $|\xi|=\delta$, vale
\[
|H_u(\zeta)|=|a e^{i\alpha\zeta}|\,|1-e^{iw\xi}|\geq c_\delta>0.
\]
La costante è uniforme: $a,w$ hanno minimi positivi, $\alpha,\gamma$
sono limitate e il secondo fattore non si annulla sul compatto indicato.
Tutte le circonferenze stanno in una striscia fissata, mentre le loro
parti reali tendono uniformemente a infinito con $m$.
Per \eqref{eq:strip}, il teorema di Rouché applicato a
$\zeta^2\mathcal F_K(\zeta u)$ dà esattamente uno zero, contato con
molteplicità, nel disco $B(\lambda_m(u),\delta)$. Esso è quindi semplice;
il disco non contiene $0$ quando $m$ è grande.

La funzione $(u,\zeta)\mapsto\mathcal F_K(\zeta u)$ è reale analitica
in $u$, in ogni carta analitica della sfera, e olomorfa in $\zeta$.
Il teorema della funzione implicita, insieme alla semplicità dello zero,
fornisce una scelta reale analitica locale. La continuità del centro
$\lambda_m(u)$ e l'unicità dello zero nel disco identificano queste
scelte nei dintorni e permettono di incollarle sulla sfera.
Ripetendo Rouché con ogni raggio $0<\delta'<\delta$ si ottiene, per $m$
sufficientemente grande, $|z_{m,K}-\lambda_m|<\delta'$ uniformemente;
questo prova il resto $o(1)$.
\end{proof}

\begin{theorem}[Unicità sotto le condizioni sulle sezioni]\label{thm:sectionunique}
Siano $K,L\in C^2_+$ in $\R^3$. Supponiamo (U) per $K$ e (S) per $L$.
Allora
\[
g_K=g_L\quad\Longrightarrow\quad L=v+K\quad\text{oppure}\quad L=v-K.
\]
\end{theorem}
\begin{proof}
L'uguaglianza dei supporti dei covariogrammi dà $w_K=w_L=w$.
L'identità di convoluzione e la trasformata intera danno, per ogni
$z\in\mathbb C^3$,
\[
\mathcal F_K(z)\mathcal F_K(-z)=\mathcal F_L(z)\mathcal F_L(-z).
\]
Per ciascuno dei rami del Lemma \ref{lem:branches}, le funzioni
\[
P_m(u)=\mathcal F_L(z_{m,K}(u)u),\qquad
Q_m(u)=\mathcal F_L(-z_{m,K}(u)u)
\]
sono reali analitiche sulla sfera connessa e soddisfano $P_mQ_m=0$.
Quindi $P_m\equiv0$ oppure $Q_m\equiv0$. Infatti, se $P_m$ non è
identicamente nulla, è non nulla su un aperto; lì $Q_m=0$, e l'identità
analitica forza $Q_m=0$ ovunque. Questo passaggio evita di costruire o
classificare i rami di zeri di $L$.

Almeno una delle due alternative si verifica lungo una successione
infinita di indici. Se $P_m\equiv0$ lungo tale successione, fissato $u$
applichiamo \eqref{eq:strip} per $L$ in $\zeta=z_{m,K}(u)$ e dividiamo
per $e^{i\alpha_L(u)\zeta}$, il cui modulo è limitato inferiormente.
Poiché
\[
e^{iw(u)z_{m,K}(u)}=-e^{-w(u)\gamma_K(u)}+o(1),
\]
il limite dà $a_L(u)=b_L(u)e^{-w(u)\gamma_K(u)}$, cioè
$\gamma_L(u)=\gamma_K(u)$ per ogni $u$.
Se invece $Q_m\equiv0$ lungo una successione infinita, la realtà
dell'indicatrice dà
\[
\mathcal F_L(-\zeta u)=\overline{\mathcal F_L(\overline\zeta u)}.
\]
Applichiamo quindi lo stesso ragionamento agli zeri
$\overline{z_{m,K}(u)}$ e otteniamo $\gamma_L=-\gamma_K$ su tutta la sfera.

Da \eqref{eq:slopes} segue
$2w\gamma_K=\log\kappa_K(-u)-\log\kappa_K(u)$, e analogamente per $L$.
Il covariogramma dà inoltre le coppie non ordinate di curvature,
come ricordato nella prima sezione. Rapporto ordinato e coppia non
ordinata impongono dunque $\kappa_L(u)=\kappa_K(u)$ per ogni $u$, oppure
$\kappa_L(u)=\kappa_K(-u)$ per ogni $u$. Nei punti con curvature
antipodali uguali le due conclusioni coincidono. L'unicità di Minkowski
conclude.
\end{proof}

\section{Verifica geometrica della condizione: Hessiane con modulo di Dini}
\label{sec:dini}
Assumiamo che i grafici del bordo nei piani tangenti abbiano, su dischi
di raggio comune, Hessiane uniformemente positive e un modulo comune
$\omega$ tale che
\begin{equation}\label{eq:dini}
\|D^2\phi(y)-D^2\phi(y')\|\leq\omega(|y-y'|),\qquad
\int_0^{r_0}\frac{\omega(r)}r\,dr<\infty.
\end{equation}
Il modulo può essere preso non decrescente. Questa formulazione rende
esplicita l'uniformità rispetto al punto di tangenza. In particolare la
soddisfa ogni corpo $C^{2,\eta}_+$, $0<\eta\leq1$: per compattezza del
bordo e cambi di coordinate regolari si può prendere $\omega(r)=Cr^\eta$.

\begin{proposition}\label{prop:dini}
Un corpo $C^2_+$ che soddisfa \eqref{eq:dini} soddisfa (U).
\end{proposition}
\begin{proof}
Vicino al punto di minimo nella direzione $u$, trasliamo il punto
all'origine e scriviamo il bordo $t=\phi(y)$, con
$\phi(0)=0$, $\nabla\phi(0)=0$ e $D^2\phi\geq\lambda I>0$.
Per $s=t-\alpha_K(u)>0$ piccolo, tutta la sezione è la sottolivellare
$\{y:\phi(y)\leq s\}$. La dimensione tangenziale è due: scrivendo
$y=r\theta$ con $\theta\in\mathbb S^1$ e
$\phi(r(s,\theta)\theta)=s$, si ottiene
\begin{align*}
A(s)&=\frac12\int_{\mathbb S^1}r(s,\theta)^2\,d\theta,\\
A'(s)&=\int_{\mathbb S^1}\frac r{\phi_r}\,d\theta,\\
A''(s)&=\int_{\mathbb S^1}
\frac{\phi_r-r\phi_{rr}}{\phi_r^3}\,d\theta.
\end{align*}
Non si è differenziato il bordo tre volte. Dalle Hessiane uniformemente
ellittiche, $r\asymp\sqrt{s}$ e $\phi_r\geq\lambda r$.
Inoltre
\[
\phi_r(r\theta)-r\phi_{rr}(r\theta)
=\int_0^r\bigl(\phi_{rr}(q\theta)-\phi_{rr}(r\theta)\bigr)\,dq,
\]
il cui modulo è al più $r\omega(r)$. Ne segue la stima uniforme
\begin{equation}\label{eq:sectiontail}
|A''(s)|\leq C\frac{\omega(C\sqrt{s})}{s}.
\end{equation}
Il membro destro è integrabile a $0$ grazie a \eqref{eq:dini}.
Inoltre, se $H=D^2\phi(0)$,
\[
\lim_{s\downarrow0}A'(s)
=\int_{\mathbb S^1}\frac{d\theta}{\theta^TH\theta}
=\frac{2\pi}{\sqrt{\det H}}.
\]
L'ultima identità si ottiene diagonalizzando la matrice positiva $H$.
Al punto di tangenza $\det H=\kappa_K(-u)$.
Lo stesso calcolo all'altro estremo dà \eqref{eq:slopes} e una stima
analoga a \eqref{eq:sectiontail}.

All'interno dell'intervallo delle sezioni, $A''(u,t)$ è congiuntamente
continua. Per verificarlo localmente, scegliamo un punto interno della
sezione, un riferimento del piano dipendente regolarmente da $u$ e una
funzione definente $C^2$ del bordo. I raggi della sezione sono funzioni
$C^2$ dei parametri per il teorema implicito: la direzione radiale è
trasversale al bordo, poiché parte da un punto interno di una sezione
convessa. La formula polare dell'area consente due derivate e dimostra
la continuità asserita. Queste osservazioni e le code integrabili mostrano
che $A'$ è assolutamente continua fino agli estremi, cioè (S).

Infine sia $u_j\to u$. Gli estremi delle sezioni convergono. Su ogni
intervallo strettamente interno la continuità appena provata dà la
convergenza delle derivate seconde in $L^1$. L'integrale di
\eqref{eq:sectiontail} su una coda di lunghezza $\delta$ tende a zero
uniformemente nelle direzioni quando $\delta\downarrow0$; lo stesso
vale per la coda superiore. Dopo il prolungamento per zero, queste
stime danno $\|T_{K,u_j}-T_{K,u}\|_1\to0$. La sfera è compatta, quindi
la sua immagine in $L^1$ è compatta: vale (U).
\end{proof}

\begin{corollary}[Conclusione nella classe con modulo di Dini]\label{cor:dini}
Siano $K,L\in C^2_+$ entrambi con l'ipotesi uniforme
\eqref{eq:dini}. Allora $g_K=g_L$ implica $L=v+K$ oppure $L=v-K$.
In particolare questa conclusione vale fra corpi $C^{2,\eta}_+$,
per ogni $0<\eta\leq1$.
\end{corollary}
\begin{proof}
Applichiamo la Proposizione \ref{prop:dini} a entrambi i corpi e quindi
il Teorema \ref{thm:sectionunique}.
\end{proof}

\paragraph{Portata e controllo ancora necessario.}
Il corollario è una deduzione con dimostrazione proposta qui, da
controllare indipendentemente anche alla luce della letteratura.
Non conclude il caso di un concorrente convesso arbitrario: l'ipotesi
sulle sezioni di $L$ è effettivamente usata in \eqref{eq:strip}.
La sola continuità delle Hessiane non rende integrabile
$\omega(\sqrt{s})/s$. Né un'approssimazione liscia di $K,L$ conserva
automaticamente l'uguaglianza dei loro covariogrammi.

\section{Trasportare l'informazione a un concorrente meno regolare}
\label{sec:transport}
Il Teorema \ref{thm:sectionunique} usa (S) anche per $L$. Il seguente
sviluppo tenta di rimuovere questa ipotesi mediante la fattorizzazione
di Hadamard e un limite sull'asse immaginario. Tutti i passaggi sono
esplicitati per consentire una revisione indipendente; la portata resta
limitata all'ipotesi (U) sul primo corpo.

\begin{lemma}[Localizzazione di tutti gli zeri scalari]\label{lem:allzeros}
Se $K\in C^2_+$ soddisfa (U), esiste $M$ tale che tutti gli zeri di
$\zeta\mapsto\mathcal F_K(\zeta u)$, per ogni $u$, appartengono alla
striscia $|\operatorname{Im}\zeta|\leq M$. Inoltre tutti gli zeri con
parte reale positiva sufficientemente grande sono esattamente i rami
$z_{m,K}(u)$ del Lemma \ref{lem:branches}.
\end{lemma}
\begin{proof}
Per $\operatorname{Im}\zeta=\tau>0$, dividiamo \eqref{eq:twoparts} per
$e^{i\alpha\zeta}$. Il termine $a$ è limitato inferiormente, mentre
\[
|b e^{iw\zeta}|\leq b e^{-w\tau},\qquad
\left|\int T_{K,u}(t)e^{i(t-\alpha)\zeta}\,dt\right|
\leq\int_\alpha^\beta |T_{K,u}(t)|e^{-\tau(t-\alpha)}\,dt.
\]
Quest'ultimo integrale tende uniformemente a zero: su
$[\alpha,\alpha+\delta]$ si usa l'uniforme integrabilità della famiglia
compatta in $L^1$, e sul resto il fattore $e^{-\tau\delta}$ e il limite
uniforme delle norme $L^1$. Lo stesso ragionamento all'altro estremo
esclude gli zeri con parte immaginaria molto negativa.

Nella striscia restante il modello $H_u$ è uniformemente separato da
zero fuori dai dischi $B(\lambda_m(u),\delta)$. Infatti, dopo aver tolto
il fattore $e^{i\alpha\zeta}$, l'espressione $a+b e^{iw\zeta}$ dipende
dalla parte reale solo attraverso una fase periodica; la separazione
dai suoi zeri dà un minimo positivo su un compatto di parametri.
L'asintotica uniforme \eqref{eq:strip} esclude gli zeri fuori dai dischi
per parte reale grande. Rouché ne aveva già contato uno in ogni disco.
\end{proof}

\paragraph{Ipotesi spettrale (Z).}
Per separare il caso $C^9_+$ dal criterio sulle sezioni, diremo che
$K\in C^2_+$ soddisfa (Z) se esiste $m_0$ e, per ogni intero
$m\geq m_0$, un ramo reale analitico $z_{m,K}:\Sph\to\mathbb C$
con le seguenti proprietà:
\begin{enumerate}
\item Per ogni $u$, gli $z_{m,K}(u)$ sono zeri semplici, distinti,
      con parte reale positiva, di $\zeta\mapsto\mathcal F_K(\zeta u)$.
\item Vale \eqref{eq:rootasymptotic}, con resto $o(1)$ uniforme in $u$.
\item Tutti gli altri zeri scalari con parte reale positiva
      appartengono a un insieme limitato comune a tutte le direzioni;
      per ogni direzione sono in numero finito, contati con molteplicità.
\end{enumerate}
La simmetria $\mathcal F_K(-\overline\zeta u)=
\overline{\mathcal F_K(\zeta u)}$ e la positività sull'asse immaginario
completano la descrizione degli zeri. In particolare essi sono tutti
in una striscia orizzontale comune. I Lemmi \ref{lem:branches} e
\ref{lem:allzeros} mostrano che (U) implica (Z). Per $K\in C^9_+$,
(Z) segue direttamente da \cite[Teoremi 4.1--4.2]{B16}, compresa la
semplicità verificata nella dimostrazione del Teorema 4.2.

\begin{lemma}[Un solo parametro per le pendenze antipodali]\label{lem:mean}
Sia $K\in C^2_+$ non centralmente simmetrico e con (Z). Sia $L$ un
corpo convesso con $g_L=g_K$. Supponiamo che, per quasi ogni $u$, le
aree di sezione di $L$ soddisfino
\begin{equation}\label{eq:linearends}
A_L(u,\alpha_L+s)=a_L(u)s+o(s),\qquad
A_L(u,\beta_L-s)=b_L(u)s+o(s),\qquad a_L,b_L>0.
\end{equation}
Allora esiste un unico numero $c\in[-1,1]$ tale che, quasi ovunque,
\begin{equation}\label{eq:scalartransport}
\frac{\log(b_L(u)/a_L(u))}{w(u)}=c\,\gamma_K(u).
\end{equation}
\end{lemma}
\begin{proof}
La funzione continua $\gamma_K$ non è identicamente nulla: altrimenti
la misura di area di $K$ sarebbe pari e l'unicità di Minkowski darebbe
simmetria centrale. Scegliamo una direzione dove $\gamma_K\ne0$.
Per tutti gli indici grandi i corrispondenti zeri di $K$ sono non reali,
semplici e distinti dai loro coniugati, per (Z). Nessun loro coniugato
è un altro zero di $K$ in questa direzione: i rami grandi sono
separati nelle parti reali, e quelli eccezionali sono limitati.
Possiamo aumentare $m_0$ affinché ciò valga per ogni $m\geq m_0$:
i rami eliminati sono in numero finito e restano limitati sulla sfera.

Da $g_K=g_L$ seguono $w_K=w_L$ e l'identità intera
$\mathcal F_K(z)\mathcal F_K(-z)=\mathcal F_L(z)\mathcal F_L(-z)$.
Per ogni indice grande le due funzioni reali analitiche
\[
P_m(u)=\mathcal F_L(z_{m,K}(u)u),\qquad
Q_m(u)=\mathcal F_L(-z_{m,K}(u)u)
\]
soddisfano $P_mQ_m=0$ sulla sfera connessa. Se $P_m$ non è
identicamente nulla, è non nulla su un aperto e lì $Q_m=0$;
l'identità analitica implica allora $Q_m\equiv0$. Quindi almeno una
è identicamente nulla. Non possono esserlo entrambe: nella direzione
appena scelta si otterrebbe una molteplicità almeno due per uno zero
semplice di $\mathcal F_K(\zeta u)\mathcal F_K(-\zeta u)$.
Esiste perciò una sequenza $\epsilon_m\in\{-1,1\}$, indipendente da
$u$, che seleziona per $L$ rispettivamente $z_{m,K}$ o il suo coniugato.
La classificazione che ne segue vale anche nelle direzioni in cui
$\gamma_K=0$. Infatti, aumentando ancora $m_0$, l'asintotica uniforme
in (Z) separa, per ogni $u$, i dischi contenenti indici diversi,
compresi i loro coniugati, e li separa dagli zeri eccezionali.
Se $z_{m,K}(u)$ non è reale, il suo coniugato non è uno zero di $K$:
il solo ramo che potrebbe contenerlo è quello di indice $m$.
Lo zero del prodotto è dunque semplice, e il fattore selezionato
esaurisce la sua molteplicità. Se è reale, la simmetria degli zeri
di $\mathcal F_L$ rispetto all'asse immaginario forza la stessa
molteplicità nei due fattori, quindi una in ciascuno.

Per ciascuna direzione fissata, tutti gli zeri di $\mathcal F_L(\zeta u)$
con parte reale positiva, tranne un numero finito, sono dunque
\begin{equation}\label{eq:chosenroots}
\zeta_{m,L}=x_m+i y_m,\qquad
x_m=\frac{(2m+1)\pi}{w}+o(1),\qquad
y_m=\epsilon_m\gamma_K(u)+o(1).
\end{equation}
Le molteplicità sono uno. Se uno zero di $K$ è reale, la molteplicità
due nel prodotto impone ancora molteplicità uno nel fattore di $L$.
Gli zeri nella metà sinistra sono $-\overline{\zeta_{m,L}}$, per la
realtà dell'indicatrice. Non vi sono zeri sull'asse immaginario, perché
$\mathcal F_L(i\tau u)>0$ per $\tau$ reale.

Fissiamo ora una direzione che soddisfi \eqref{eq:linearends} e
scriviamo $F(\zeta)=\mathcal F_L(\zeta u)$.
Raggruppiamo tutti gli zeri, compresi quelli eccezionali e contati con
molteplicità, nelle coppie $\zeta_j=x_j+iy_j,-\overline{\zeta_j}$,
con $x_j>0$. Abbiamo $|y_j|\leq M$, $x_j\asymp j$ per $j$ grande.
In particolare
\[
\sum_j\left|\frac1{\zeta_j}-\frac1{\overline{\zeta_j}}\right|<\infty,
\qquad \sum_j|\zeta_j|^{-2}<\infty.
\]
La fattorizzazione di Hadamard, raggruppando i fattori canonici di
genere uno e assorbendone gli esponenziali lineari, diventa
\begin{equation}\label{eq:pairedproduct}
F(\zeta)=F(0)e^{iq\zeta}
\prod_j\left(1-\frac{\zeta}{\zeta_j}\right)
       \left(1+\frac{\zeta}{\overline{\zeta_j}}\right),\qquad q\in\R.
\end{equation}
Per controllare l'ipotesi di crescita del teorema di Hadamard,
$L\subset B(0,T)$ implica
$|F(\zeta)|\leq\vol(L)e^{T|\operatorname{Im}\zeta|}$.
Il polinomio nell'esponenziale ha quindi grado al più uno; si veda
anche \cite[p.~4]{Siu24}. Non rimane un fattore intero arbitrario.
Il prodotto converge uniformemente sui compatti: il fattore doppio
meno uno è $O((|\zeta|+|\zeta|^2)/x_j^2)$ su ogni compatto.
Il coefficiente lineare nell'esponenziale è immaginario puro perché sia
$F'(0)/F(0)$ sia la derivata logaritmica del prodotto in zero sono
immaginarie pure. Qui si usa $F(0)=\vol(L)>0$.

Per $\tau>0$ poniamo
\[
D_L(\tau)=\log F(i\tau)-\log F(-i\tau).
\]
Tutti i fattori doppi in \eqref{eq:pairedproduct} sono positivi
sull'asse immaginario. Possiamo quindi usare logaritmi reali e ottenere
\begin{equation}\label{eq:logproduct}
D_L(\tau)=-2q\tau+\sum_j R_{x_j,y_j}(\tau),\qquad
R_{x,y}(\tau)=
\log\frac{x^2+(\tau-y)^2}{x^2+(\tau+y)^2}.
\end{equation}
La serie converge assolutamente per ogni $\tau$ fissato.
Per $|y|\leq M$ e $\tau\geq2M+1$,
$|R_{x,y}(\tau)|\leq C_M\tau/(x^2+\tau^2)$.
Poiché gli $x_j$ crescono linearmente, la serie in
\eqref{eq:logproduct} è limitata uniformemente per $\tau\to\infty$.

D'altra parte, \eqref{eq:linearends} dà direttamente
\begin{equation}\label{eq:laplaceends}
F(i\tau)=e^{-\alpha_L\tau}\frac{a_L}{\tau^2}(1+o(1)),\qquad
F(-i\tau)=e^{\beta_L\tau}\frac{b_L}{\tau^2}(1+o(1)).
\end{equation}
Per esempio, dopo aver posto $s=t-\alpha_L$, il cambio di scala
$r=\tau s$ e $A_L(\alpha_L+s)/s\to a_L$ danno il primo limite;
il contributo lontano dall'estremo è esponenzialmente piccolo.
Dividendo \eqref{eq:logproduct} per $\tau$ e confrontando con
\eqref{eq:laplaceends}, si ricava $2q=\alpha_L+\beta_L$.
Pertanto la quantità centrata
\begin{equation}\label{eq:centeredlog}
B_L(\tau):=D_L(\tau)+(\alpha_L+\beta_L)\tau
=\sum_j R_{x_j,y_j}(\tau)
\end{equation}
soddisfa
\begin{equation}\label{eq:laplaceJ}
\lim_{\tau\to\infty}B_L(\tau)=\log(a_L/b_L).
\end{equation}
Questo argomento non differenzia alcun resto asintotico né la serie.

Lo sviluppo $\log[(1-v)/(1+v)]=-2v+O(v^3)$, con
$v=2\tau y/(x^2+\tau^2+y^2)$, dà uniformemente in $x>0$
\begin{equation}\label{eq:kernel}
R_{x,y}(\tau)=-\frac{4y\tau}{x^2+\tau^2}
+O_M\left(\frac{\tau}{(x^2+\tau^2)^2}\right).
\end{equation}
Infatti $|v|$ resta separato da uno; sia la correzione del denominatore
sia il termine cubico sono maggiorati dal resto scritto.
Più precisamente $|v|\leq4/5$ e
\[
\left|\log\frac{1-v}{1+v}+2v\right|
\leq\frac{2|v|^3}{3(1-|v|^2)}.
\]
Posto $A=x^2+\tau^2$, la correzione fra $-2v$ e $-4y\tau/A$
è al più $4M^3\tau/A^2$. Il resto complessivo è dunque al più
$20M^3\tau/A^2$. Poiché $x_j\asymp j$, il confronto con
$\int_0^\infty\tau/(\rho^2t^2+\tau^2)^2\,dt=\pi/(4\rho\tau^2)$, $\rho>0$,
giustifica la sommabilità uniforme del resto per $\tau\to\infty$.
La somma dei resti è $O(\tau^{-2})$, mentre ogni coppia eccezionale
fissata contribuisce $o(1)$. Sostituendo \eqref{eq:chosenroots} si trova
\[
B_L(\tau)=-4\gamma_K(u)\sum_{m\geq m_0}
\frac{\epsilon_m\tau}{((2m+1)\pi/w)^2+\tau^2}+o(1).
\]
Per gli $y_m$ si separa una testa finita da una coda con errore
uniformemente piccolo e si usa la somma limitata dei pesi positivi.
Per gli $x_m$, il loro scarto dal modello è limitato e la somma dei
moduli delle derivate di $\tau/(x^2+\tau^2)$ è $O(\tau^{-1})$;
la sostituzione cambia dunque la somma di $o(1)$.
Ponendo $\tau=2\pi R/w$ otteniamo
\begin{equation}\label{eq:universalmean}
B_L(2\pi R/w)=-w\gamma_K(u)M(R)+o(1),\qquad
M(R)=\frac2\pi\sum_{m\geq m_0}
\frac{\epsilon_m R}{(m+1/2)^2+R^2}.
\end{equation}
La media $M(R)$ non dipende da $u$ e ha pesi positivi. La loro somma
tende a $\frac2\pi\int_0^\infty(1+t^2)^{-1}\,dt=1$, quindi
$\limsup|M(R)|\leq1$.

Esiste una direzione con \eqref{eq:linearends} e $\gamma_K\ne0$, perché
quest'ultima disuguaglianza vale su un aperto. In tale direzione
\eqref{eq:universalmean}--\eqref{eq:laplaceJ} provano che $M(R)$ ha un
limite $c\in[-1,1]$. Riutilizzando il medesimo limite in ogni direzione
ammissibile otteniamo \eqref{eq:scalartransport}.
\end{proof}

\begin{remark}[Quali ipotesi usa il trasporto]
Il Lemma \ref{lem:mean} usa (Z), non (U). Per il caso $C^9_+$,
l'ipotesi spettrale è fornita dal risultato pubblicato di Bianchi;
per il criterio esteso sulle sezioni, è fornita dai Lemmi
\ref{lem:branches} e \ref{lem:allzeros}. Sul concorrente $L$ si usano
soltanto le espansioni \eqref{eq:linearends}, l'identità delle
trasformate e le proprietà di una funzione intera di tipo esponenziale.
\end{remark}

\paragraph{Seconda derivazione del limite, mediante differenze finite.}
Si può verificare il passaggio spettrale senza identificare il
coefficiente $q$ dell'esponenziale in \eqref{eq:pairedproduct}.
Con la stessa $D_L$, poniamo
\[
H_L(\tau)=D_L(2\tau)-2D_L(\tau).
\]
Ogni termine lineare in $\tau$ si cancella esattamente. La
fattorizzazione di Hadamard dà quindi
\[
H_L(\tau)=\sum_j\bigl(R_{x_j,y_j}(2\tau)-2R_{x_j,y_j}(\tau)\bigr),
\]
mentre le espansioni di Laplace \eqref{eq:laplaceends} danno
$H_L(\tau)\to\log(b_L/a_L)$.
Per $|y|\leq M$ e $\tau\geq2M+1$, \eqref{eq:kernel} implica
\[
R_{x,y}(2\tau)-2R_{x,y}(\tau)
=\frac{24y\tau^3}{(x^2+\tau^2)(x^2+4\tau^2)}
+O_M\!\left(\frac{\tau}{(x^2+\tau^2)^2}\right).
\]
La somma dei resti tende a zero. Con le stesse sostituzioni degli
zeri e $\tau=2\pi R/w$ si ottiene
\[
H_L(2\pi R/w)=w\gamma_K(u)\,N(R)+o(1),\qquad
N(R)=\frac{12}{\pi}\sum_{m\geq m_0}
\frac{\epsilon_m R^3}
{((m+1/2)^2+R^2)((m+1/2)^2+4R^2)}.
\]
Anche questi pesi sono positivi e la loro somma tende a uno, poiché
\[
\frac{12}{\pi}\int_0^\infty
\frac{dt}{(1+t^2)(4+t^2)}=1.
\]
Una direzione ammissibile con $\gamma_K\ne0$ forza l'esistenza del
limite di $N(R)$; lo stesso limite in ogni direzione ridà
\eqref{eq:scalartransport}. È una seconda derivazione del passaggio
al limite, con un nucleo diverso, che elimina ogni dipendenza dal
valore di $q$. Non sostituisce le verifiche dei rami di zeri e del
lemma geometrico che segue.

\subsection*{Dati quadratici quasi ovunque del concorrente}
Per il prossimo passaggio usiamo il teorema di Rademacher e la formula
dell'area per mappe Lipschitz. Dopo aver dedotto il supporto $C^{1,1}$,
Rademacher applicato al suo gradiente fornisce le espansioni quadratiche
quasi ovunque. Esplicitiamo l'argomento per il concorrente arbitrario.

\begin{lemma}\label{lem:roughcompetitor}
Sia $K\in C^2_+$ e sia $L$ convesso con $g_L=g_K$. La funzione di
supporto omogenea di $L$ è $C^{1,1}$ lontano dall'origine. Quasi ovunque
su $\Sph$ esistono matrici positive definite $R_L(u)$, con
\[
dS_L(u)=\det R_L(u)\,d\sigma(u),\qquad
\{\det R_L(u),\det R_L(-u)\}=
\{f_K(u),f_K(-u)\}.
\]
In queste direzioni valgono \eqref{eq:linearends}, con
$a_L=2\pi\sqrt{\det R_L(-u)}$ e
$b_L=2\pi\sqrt{\det R_L(u)}$.
\end{lemma}
\begin{proof}
Poniamo $D=K-K=L-L$. Le Hessiane distribuzionali delle funzioni di
supporto omogenee soddisfano, come misure matriciali,
\[
0\leq D^2h_L\leq D^2h_D,
\]
poiché $h_D=h_L+h_{-L}$. Su ogni anello lontano dall'origine il membro
destro ha una densità continua limitata. Dunque $h_L\in C^{1,1}$ e
la mappa $p_L(u)=\nabla h_L(u)$ è Lipschitz. Essa parametrizza tutto il
bordo, e il suo Jacobiano tangenziale è $\det R_L(u)$, dove
$R_L=\nabla_{\Sph}^2h_L+h_LI\succeq0$ quasi ovunque.
La formula dell'area dà $dS_L=\det R_L\,d\sigma$: nei punti regolari
del bordo la normale è unica; i punti non regolari hanno area nulla.
Più precisamente, per ogni insieme boreliano $E\subset\Sph$, la
molteplicità della restrizione $p_L|_E$ in un punto regolare $x$ vale
$\ind_E(\nu_L(x))$. La formula dell'area, integrata su questi punti,
dà $\int_E\det R_L\,d\sigma=S_L(E)$; l'insieme non regolare non
contribuisce. Questo esclude anche una parte singolare della misura
di area, senza aver assunto un bordo $C^2$ per $L$.

Fissiamo una direzione in cui esistano le derivate seconde del
supporto in $u$ e in $-u$; queste direzioni hanno misura piena.
Infatti, nelle carte $v\mapsto(v,1)$ il gradiente del supporto è
Lipschitz e quindi differenziabile quasi ovunque. Integrando la
sua espansione del primo ordine lungo un segmento si ottiene
il corrispondente sviluppo di Taylor quadratico.
Studiamo dapprima soltanto l'estremo superiore, scegliendo $u=e_3$
e $p_L(u)=0$. Nel riferimento tangente, con $R=R_L(u)\succeq0$,
\[
G(v):=h_L(v,1)=\tfrac12 v^T Rv+o(|v|^2),\qquad v\in\R^2.
\]
Mostriamo, includendo il caso degenere, che
\begin{equation}\label{eq:allslopes}
\lim_{s\downarrow0}\frac{A_L(u,\beta_L-s)}s
=2\pi\sqrt{\det R_L(u)}.
\end{equation}

Le sezioni vicine all'estremo si restringono al punto esposto $0$,
perché il supporto è differenziabile. Dalla stima
$G(v)\leq C|v|^2$, scegliendo $v=y/(2C)$ nella disuguaglianza
$v\cdot y-s\leq G(v)$, segue $|y|^2\leq4Cs$ per $s$ piccolo.
Le sezioni riscalate $E_s=\{Y:(\sqrt{s}Y,-s)\in L\}$ sono dunque
uniformemente limitate. Per ogni loro punto limite $Y$, ponendo
$v=\sqrt{s}V$ nelle disuguaglianze di supporto otteniamo
\[
V\cdot Y-1\leq\tfrac12 V^T RV\qquad\text{per ogni }V\in\R^2.
\]
Se $R$ è singolare, scegliendo $V$ arbitrario nel suo nucleo si
ricava $Y\in\operatorname{im}R$. Questo sottospazio ha dimensione
al più uno. Per compattezza, gli $E_s$ sono infine contenuti in ogni
intorno prefissato di questo sottospazio entro un disco limitato.
La loro area tende quindi a zero, provando \eqref{eq:allslopes}
quando $\det R=0$.

Se invece $R>0$, il punto $p_L(u)$ ha una sola normale. Infatti, una
seconda normale produrrebbe, per convessità del cono normale, un
segmento uscente da zero nella carta $(v,1)$ sul quale $G(v)=0$;
ciò contraddice l'espansione quadratica positiva. Il punto $y=0$ è
interno alla proiezione di $L$ sul piano tangente: altrimenti un
piano di supporto verticale darebbe un'altra normale. Il bordo
superiore si scrive quindi localmente $z=-\phi(y)$, con $\phi$
convessa e finita su un disco e $\partial\phi(0)=\{0\}$.

La dualità di Legendre fra $G$ e $\phi$ è locale in entrambi i sensi.
L'unicità del punto esposto implica che i punti di contatto tendano
a zero quando $v\to0$. Viceversa, i sottogradienti di $\phi$ sono
localmente limitati e ogni loro limite per $y\to0$ appartiene a
$\partial\phi(0)=\{0\}$. Per ogni $\varepsilon>0$ piccolo,
le stime locali
\[
\tfrac12 v^T(R-\varepsilon I)v\leq G(v)
\leq\tfrac12 v^T(R+\varepsilon I)v
\]
si trasformano quindi nelle stime
\[
\tfrac12 y^T(R+\varepsilon I)^{-1}y\leq\phi(y)
\leq\tfrac12 y^T(R-\varepsilon I)^{-1}y.
\]
I sottolivelli $\{\phi\leq s\}$, che descrivono le sezioni vicine
all'estremo, sono racchiusi fra le corrispondenti ellissi.
Per giustificare che il bordo inferiore della fibra non interferisce,
osserviamo che sopra un punto interno alla proiezione la fibra di un
corpo convesso pieno ha lunghezza positiva. Il suo estremo inferiore
è una funzione convessa, finita e continua nell'interno della
proiezione. In un intorno di $y=0$ esso resta quindi sotto $-d$ per
qualche $d>0$. Le sezioni si restringono a zero, perciò per $s$ piccolo
la sola condizione effettiva è $\phi(y)\leq s$. Dividendo
le aree per $s$ e poi mandando $\varepsilon$ a zero si ottiene
\eqref{eq:allslopes}, perché l'area di
$\{Y:Y^TR^{-1}Y/2\leq1\}$ è $2\pi\sqrt{\det R}$.

Lo stesso argomento in $-u$ fornisce le due espansioni di
\eqref{eq:linearends} con
coefficienti inizialmente soltanto non negativi e finiti,
\[
a_L=2\pi\sqrt{\det R_L(-u)},\qquad
b_L=2\pi\sqrt{\det R_L(u)}.
\]
Non abbiamo ancora usato la loro positività. Essa è forzata proprio
dal covariogramma: l'identità
$\mathcal F_L(z)\mathcal F_L(-z)=\mathcal F_K(z)\mathcal F_K(-z)$,
valutata in $z=i\tau u$, e l'uguaglianza delle larghezze danno
\begin{equation}\label{eq:laplaceproduct}
a_Lb_L
=\lim_{\tau\to\infty}\tau^4e^{-w\tau}
  \mathcal F_L(i\tau u)\mathcal F_L(-i\tau u)
=a_Kb_K>0.
\end{equation}
Il limite di Laplace vale anche quando un coefficiente è zero:
$A(\alpha+s)=o(s)$ implica
$\tau^2 e^{\alpha\tau}\mathcal F(i\tau u)\to0$.
Dunque entrambi i determinanti di $L$ sono positivi e
\[
\det R_L(u)\det R_L(-u)=f_K(u)f_K(-u).
\]
Questo recupera il prodotto senza calcolare un'intersezione di calotte
né postulare un bordo regolare per il concorrente.

Infine le derivate direzionali del covariogramma all'origine danno le
aree delle proiezioni. L'iniettività della trasformata coseno sulle
misure pari implica $S_L+S_{-L}=S_K+S_{-K}$, dunque l'uguaglianza
delle somme dei determinanti quasi ovunque. Somma e prodotto danno
la coppia non ordinata dell'enunciato.
\end{proof}

\begin{theorem}[Estensione proposta ai concorrenti arbitrari]\label{thm:allcompetitors}
Sia $K\in C^2_+$ in $\R^3$ e supponiamo (U) per $K$.
Per ogni corpo convesso $L$ in $\R^3$, senza un'ipotesi di regolarità
su $L$, vale
\[
g_L=g_K\quad\Longrightarrow\quad L=v+K\quad\text{oppure}\quad L=v-K.
\]
In particolare, se i passaggi della presente sezione sono confermati
dalla revisione indipendente, il criterio di Dini della Sezione
\ref{sec:dini} determina il corpo anche fra concorrenti arbitrari.
\end{theorem}
\begin{proof}
Se $K$ è centralmente simmetrico, $\vol(D)=8\vol(K)=8\vol(L)$.
L'uguaglianza in Brunn--Minkowski per $L+(-L)$ impone la simmetria
centrale di $L$; $D$ determina quindi entrambi a meno di traslazione.

Altrimenti (U) implica (Z), e il Lemma \ref{lem:roughcompetitor} dà
\eqref{eq:linearends} quasi ovunque e inoltre
\[
\gamma_L(u):=\frac{\log(b_L(u)/a_L(u))}{w(u)}
\in\{\gamma_K(u),-\gamma_K(u)\}\quad\text{quasi ovunque}.
\]
Il Lemma \ref{lem:mean} fornisce nello stesso tempo
$\gamma_L=c\gamma_K$ con un solo numero $c\in[-1,1]$.
Su un aperto $\gamma_K\ne0$; scegliendovi una direzione ammissibile,
le due relazioni impongono $c=1$ oppure $c=-1$.
Rapporto e coppia non ordinata danno allora
$\det R_L(u)=f_K(u)$ quasi ovunque, oppure
$\det R_L(u)=f_K(-u)$ quasi ovunque. Le misure di area coincidono
quindi con $S_K$ oppure $S_{-K}$. L'unicità di Minkowski conclude.
\end{proof}

\begin{corollary}[Caso $C^9_+$ del programma proposto]\label{cor:targetC9}
La catena di argomenti della presente sezione dà precisamente
\eqref{eq:target}: un corpo $K\in C^9_+$ in $\R^3$ è determinato
dal proprio covariogramma fra tutti i corpi convessi, a meno di
traslazione e riflessione centrale.
\end{corollary}
\begin{proof}
Per $K\in C^9_+$, \cite[Teoremi 4.1--4.2 e dimostrazione del
Teorema 4.2, pp.~14--15 del manoscritto]{B16} fornisce già i rami
analitici globali, la semplicità e la classificazione degli zeri
scalari grandi. Per $n=3$ il termine reale della loro asintotica è
$(2m+1)\pi/w_K$, quello immaginario è $\gamma_K$ e il resto è
$O(m^{-1})$ uniforme. Gli altri zeri appartengono a un insieme
limitato. Sono esattamente le proprietà usate nel Lemma
\ref{lem:mean}, come osservato sopra.

Il Lemma \ref{lem:roughcompetitor} fornisce, per un concorrente
convesso arbitrario, le pendenze agli estremi e le coppie non ordinate
quasi ovunque. Se $K$ non è centralmente simmetrico, il Lemma
\ref{lem:mean}, applicato mediante (Z), dà un solo $c\in[-1,1]$ con
$\gamma_L=c\gamma_K$ quasi ovunque. La coppia non ordinata dà anche
$\gamma_L\in\{\gamma_K,-\gamma_K\}$ quasi ovunque. La funzione
continua $\gamma_K$ è non nulla su un aperto; scegliendovi una
direzione appartenente all'insieme di misura piena dei due lemmi
si ricava $c=1$ oppure $c=-1$.

Se $c=1$, il rapporto ordinato e la coppia non ordinata impongono
$\det R_L(u)=f_K(u)$ quasi ovunque; se $c=-1$, impongono
$\det R_L(u)=f_K(-u)$ quasi ovunque. Nei punti con coppia uguale
le due possibilità coincidono. La formula dell'area nel Lemma
\ref{lem:roughcompetitor} dà quindi $S_L=S_K$ oppure $S_L=S_{-K}$
come misure, e l'unicità di Minkowski conclude.

Se $K$ è centralmente simmetrico, poniamo $D=K-K=L-L$.
Da $g_K(0)=g_L(0)$ segue
$\vol(D)=8\vol(K)=8\vol(L)$. È il caso di uguaglianza in
Brunn--Minkowski per $L+(-L)$; quindi $L$ è centralmente simmetrico.
Il corpo differenza determina allora entrambi a meno di traslazione.

Questa dimostrazione usa (Z) da \cite{B16}, i due lemmi appena
richiamati e l'unicità di Minkowski; il criterio di Dini e l'ipotesi
(U) non intervengono nel caso $C^9_+$.
\end{proof}

\paragraph{Stato di questo sviluppo.}
La Sezione \ref{sec:transport} è il passaggio più delicato e nuovo per
questo quaderno. Non se ne rivendica la novità bibliografica e non si
sostituisce una revisione indipendente con la compilazione del testo.
La fattorizzazione canonica, il limite della media positiva e il lemma
sui dati quadratici quasi ovunque sono stati sviluppati e sottoposti
ai controlli interni descritti nel testo; non è stata individuata una
lacuna nella catena principale. Sono anche i passaggi centrali da
riesaminare indipendentemente. Nessuna regolarità di $L$ è assunta:
il supporto $C^{1,1}$ e le espansioni delle sezioni sono dedotti
dall'uguaglianza dei covariogrammi. La stesura della prova candidata
del Corollario \ref{cor:targetC9} è completa in questo senso, senza
rinviare la conclusione a un ulteriore lemma non dimostrato.
Il confronto bibliografico è rilevante: \cite[Osservazione 5.2]{B16}
segnala espressamente il limite ai concorrenti regolari del proprio
teorema; \cite[p.~23 del manoscritto, dopo il Teorema 3.25]{B23}
indica ancora esplicitamente come aperta l'estensione ai concorrenti
convessi arbitrari. Il Corollario \ref{cor:targetC9} è quindi il
bersaglio bibliografico preciso, distinto dall'affinamento $C^{2,\eta}_+$.
L'estensione proposta qui non va presentata come una
semplice citazione o come un risultato già acquisito in letteratura.
Anche se confermato integralmente, il risultato non copre un primo
corpo convesso arbitrario: il caso conclusivo richiede $K\in C^9_+$;
l'estensione accessoria del Teorema \ref{thm:allcompetitors} richiede
invece (U). La ricerca si arresta su questa versione per il riesame
richiesto dall'utente.

\section{Aggirare direzioni eccezionali senza perdere la scelta globale}
\label{sec:exceptional}
La compattezza uniforme su tutta la sfera può essere sostituita da una
condizione locale su un insieme connesso di misura piena. Questo
affinamento usa ancora il trasporto proposto nella Sezione
\ref{sec:transport}, con il medesimo stato di revisione.

\begin{theorem}[Versione locale del criterio sulle sezioni]\label{thm:localU}
Sia $K\in C^2_+$. Supponiamo che esista un aperto connesso
$\Omega\subset\Sph$ il cui complementare ha misura sferica nulla,
tale che (S) valga per $u\in\Omega$ e che
\[
u\longmapsto T_{K,u}\quad\text{sia continua da }\Omega\text{ in }L^1(\R).
\]
Allora la dimostrazione del Teorema \ref{thm:allcompetitors} si estende
a $K$: ogni corpo convesso $L$ con $g_L=g_K$ è una traslazione di
$K$ oppure di $-K$.
\end{theorem}
\begin{proof}
Il caso centralmente simmetrico è già risolto. Altrimenti scegliamo
$u_0\in\Omega$ con $\gamma_K(u_0)\ne0$ e un'esaustione crescente di
$\Omega$ mediante aperti connessi $\Omega_j$ con chiusura compatta in
$\Omega$, tutti contenenti $u_0$.
Su ciascuna chiusura la famiglia $T_{K,u}$ è compatta in $L^1$.
Le dimostrazioni dei Lemmi \ref{lem:strip}, \ref{lem:branches} e
\ref{lem:allzeros} valgono quindi uniformemente su $\Omega_j$, con
una soglia $m_j$ che può crescere con $j$.

Su $\Omega_1$ definiamo la sequenza $\epsilon_m$ mediante l'identità
analitica dei due fattori. Su $\Omega_j$, per $m$ sufficientemente
grande, il ramo coincide con quello di $\Omega_1$ nell'intersezione:
entrambi sono lo zero unico vicino allo stesso $\lambda_m(u)$.
La scelta del fattore su $\Omega_j$ coincide quindi con $\epsilon_m$,
perché $\Omega_j$ è connesso e lo zero in $u_0$ è non reale e semplice.

Per ogni $u\in\Omega$ otteniamo dunque \eqref{eq:chosenroots} con la
\emph{stessa} sequenza di segni, a meno di un numero finito di indici
dipendente da $u$. Questi indici non modificano il limite di
$B_L$ né la media $M(R)$, perché ogni peso fissato tende a zero.
La prova del Lemma \ref{lem:mean} prosegue su $\Omega$. Il Lemma
\ref{lem:roughcompetitor} fornisce i dati quadratici quasi ovunque;
poiché $\Omega$ ha misura piena, il confronto $c\gamma_K=\pm\gamma_K$
e l'uguaglianza delle misure di area concludono come prima.
\end{proof}

\begin{corollary}[Irregolarità isolate del bordo]\label{cor:isolated}
Il precedente argomento si applica a un corpo $K\in C^2_+$ il cui
bordo sia localmente $C^{2,\eta}$, con esponente positivo anche variabile
da carta a carta, fuori da un numero finito di punti.
\end{corollary}
\begin{proof}
Togliamo dalla sfera le normali dei punti eccezionali e le loro
antipodali. Il complementare $\Omega$ è aperto, connesso e di misura
piena. Su ogni compatto di $\Omega$, gli estremi delle sezioni stanno
in un numero finito di carte con un modulo di Dini comune.
Le stime di coda della Proposizione \ref{prop:dini} sono quindi
uniformi su quel compatto. Le porzioni interne delle sezioni richiedono
soltanto la regolarità $C^2$ globale. Ne segue la continuità locale
in $L^1$ richiesta dal Teorema \ref{thm:localU}.
\end{proof}

La connessione qui svolge un ruolo preciso: rende indipendente dalla
direzione il segno di ciascun ramo, quando l'indice è abbastanza grande
su un dato compatto. Non si è dimostrato che, per un corpo $C^2_+$
arbitrario, esista un aperto connesso di misura piena con tali proprietà.
In particolare non si può sostituire questa ipotesi con la sola
differenziabilità seconda quasi ovunque.

\section{Un ostacolo reale all'estensione automatica a $C^2_+$}
\label{sec:nonlone}
La condizione (S) non vale per ogni corpo $C^2_+$. Ecco una costruzione
locale che si può inserire in un corpo di rivoluzione chiuso.
Per $0<s<s_0<e^{-2}$ definiamo
\[
q'(s)=\int_{1/s}^{\infty}\frac{\sin v}{v\log v}\,dv,
\qquad q(s)=\int_0^s q'(t)\,dt.
\]
L'integrale improprio converge per Dirichlet. Integrando per parti si ha
$q'(s)=O(s/\log(1/s))$, $q(s)=O(s^2/\log(1/s))$, mentre
\[
q''(s)=\frac{\sin(1/s)}{s\log(1/s)},\qquad
s q''(s)\longrightarrow0,
\qquad\int_0^{s_0}|q''(s)|\,ds=\infty.
\]
L'ultima divergenza equivale a quella di
$\int_{1/s_0}^{\infty}|\sin v|/(v\log v)\,dv$ e segue confrontando
le parti centrali dei successivi semiperiodi con la serie
$\sum_{n\geq2}1/(n\log n)$.

Sia $\chi$ un cutoff liscio, uguale a uno per $s<s_0/3$ e nullo per
$s>2s_0/3$. Poniamo, su $[0,2]$,
\[
f(s)=2s-s^2+\delta\chi(s)q(s),\qquad
K=\{(y,s)\in\R^2\times[0,2]: |y|^2\leq f(s)\},
\]
con $\delta>0$ sufficientemente piccolo, e $q(0)=0$.
Per $\delta=0$ si ottiene la sfera unitaria traslata e
$2ff''-(f')^2=-4$. I termini perturbativi in questa espressione sono
uniformemente limitati e tendono a zero al polo inferiore, poiché
$q'=o(1)$, $q=o(s)$ e $s q''=o(1)$. Per $\delta$ piccolo rimangono
$f>0$ all'interno e $2ff''-(f')^2<0$; dunque $\sqrt f$ è strettamente
concava e il corpo di rivoluzione è convesso con curvature positive
fuori dai poli.

Al polo inferiore il grafico è $s=\phi(r)=f^{-1}(r^2)$. Si ha
\[
\frac{\phi'(r)}r=\frac2{f'(s)}\longrightarrow1,
\qquad
\phi''(r)=\frac2{f'(s)}-\frac{4r^2 f''(s)}{f'(s)^3}
\longrightarrow1.
\]
Quindi il grafico radiale è $C^2$ anche al polo, con Hessiana identità.
Vicino al polo superiore è esattamente sferico. Ne risulta $K\in C^2_+$.
Tuttavia le aree delle sezioni orizzontali sono $A(s)=\pi f(s)$ e
$A''\notin L^1$ vicino a $0$.

Questo esempio non mette in discussione l'unicità del covariogramma.
Dimostra precisamente che non si può eliminare l'ipotesi (S) affermando
che essa segue dalla sola regolarità $C^2_+$.
Il bordo dell'esempio è tuttavia liscio fuori dal polo inferiore:
rientra quindi nella versione con direzioni eccezionali del Corollario
\ref{cor:isolated}, se l'estensione della Sezione \ref{sec:transport}
è confermata. Il fallimento di (S) in una direzione non è, da solo,
un ostacolo definitivo alla strada analitica.

\section{Che cosa resta da dimostrare}
Il criterio dei momenti risolve i casi senza sottoinsiemi propri bilanciati.
Non elimina i casi bilanciati: l'esempio esplicito mostra che nemmeno
l'intera informazione locale vicino a origine e bordo del supporto
può eliminarli tutti. Occorrono traslazioni interne.

Il teorema sulla parte pari elimina tutti gli scambi parziali che preservano
$\ind_K+\ind_{-K}$ rispetto a un centro comune. Non è dimostrato che
ogni coppia di corpi convessi con lo stesso covariogramma possa essere
ridotta a questa situazione. In particolare, il corpo differenza fornisce
$h_K(u)+h_K(-u)$, non $\ind_K(x)+\ind_K(-x)$.

Una prossima domanda circoscritta è studiare scambi di porzioni riferiti
a centri diversi: in quel caso la fattorizzazione \eqref{eq:cross} con
un'unica parte pari non è disponibile. L'argomento degli zeri scarta
invece le scelte di curvatura bilanciate sotto (U); al di fuori di
questa ipotesi la loro esclusione resta da dimostrare.

Le Sezioni \ref{sec:zeros}--\ref{sec:dini} forniscono una terza strada:
rendere globale la scelta con l'analiticità della trasformata, dopo aver
costruito un ramo di zeri con asintotica uniforme. La Sezione
\ref{sec:transport} propone di trattare un concorrente arbitrario
trasportando le pendenze mediante una media positiva degli scambi.
Occorre revisionare questa estensione e, per la domanda generale,
eliminare l'ipotesi (U) anche sul primo corpo. La costruzione della
Sezione \ref{sec:nonlone} mostra perché quest'ultimo passaggio non è
automatico neppure per un corpo $C^2_+$. I bordi non lisci, le facce
e le curvature degeneri richiedono inoltre un argomento ulteriore.
La Sezione \ref{sec:exceptional} permette di aggirare eccezioni isolate,
ma non prova l'esistenza delle direzioni regolari connesse necessarie
per un corpo generico.

\paragraph{Verifiche.}
Il file \texttt{verifica\_scambi.py} controlla con aritmetica razionale le
identità polinomiali per $P_3$, i raccordi e le derivate del cutoff, il
margine positivo di convessità e le disuguaglianze per il punto testimone
e la corona. Non sostituisce le dimostrazioni e non è un proof checker.
Il file \texttt{verifica\_media\_zeri.py} confronta inoltre
$J_L(\tau)=B_L(\tau)-\tau B_L'(\tau)$, usato in una prima versione
dell'argomento, calcolato direttamente dagli integrali di Laplace con la
somma numerica sui primi $10\,000$ zeri per
$A(t)=(1-t^2)(1+t/5)$, $-1\leq t\leq1$.
Per $\tau=20,100,500$ gli scarti osservati sono rispettivamente circa
$4.4\cdot10^{-11}$, $5.5\cdot10^{-9}$ e $6.9\cdot10^{-7}$.
La versione semplificata della prova usa direttamente $B_L$ e la
somma dei logaritmi in \eqref{eq:centeredlog}. Lo stesso script ne
confronta i valori diretti con la somma troncata e con una stima
integrale della coda del modello; gli scarti dopo quest'ultima
correzione sono circa $5.3\cdot10^{-14}$, $2.6\cdot10^{-13}$ e
$1.3\cdot10^{-12}$. La correzione della coda non è un limite
d'errore certificato. Questa seconda formula evita del tutto le
derivate nella dimostrazione analitica.
La somma è troncata e gli zeri sono calcolati con Newton in virgola
mobile: è una verifica diagnostica di segni e costanti, non una
certificazione degli zeri né una dimostrazione del limite.

\section{Controllo formale parziale in Lean}\label{sec:leancheck}
Il controllo del 9 ottobre 2026 usa Lean 4.24.0 e mathlib alla
revisione \texttt{f897ebcf72cd16f89ab4577d0c826cd14afaafc7}.
Il progetto contiene 33 teoremi espliciti, distribuiti in nove moduli.
La compilazione e la rielaborazione diretta di tutti i sorgenti hanno
avuto esito positivo, trattando anche gli avvisi come errori.
L'audit delle dipendenze comprende 71 dichiarazioni di teoremi,
contando gli ausiliari generati dalle tattiche: utilizza soltanto
\texttt{propext}, \texttt{Classical.choice} e \texttt{Quot.sound},
le usuali basi classiche di Lean/mathlib. Non compaiono dipendenze
da lacune ammesse o assiomi matematici aggiunti.

\paragraph{Passaggi verificati.}
Il codice controlla il recupero di una coppia da somma e prodotto,
l'ordinamento tramite il rapporto logaritmico, il principio di
identità per la scelta di un fattore analitico su un aperto
preconnesso e l'assegnazione esclusiva di uno zero semplice di un
prodotto. Controlla inoltre il passaggio da un parametro comune
a un unico orientamento quasi ovunque e l'uguaglianza delle
misure definite dalle densità, assumendo esplicitamente le
identità di somma, prodotto e trasporto. Quest'ultima affermazione
non identifica ancora tali misure con le misure di area dei corpi.

Per il nucleo di \eqref{eq:logproduct} è verificata la stima
\[
\begin{gathered}
\left|R_{x,y}(\tau)+\frac{4\tau y}{x^2+\tau^2}\right|
\leq\frac{36M^3\tau}{(x^2+\tau^2)^2}
\leq\frac{36M^3}{\tau x^2},\\
x\ne0,\quad |y|\leq M,\quad \tau>0,\quad \tau\geq4M.
\end{gathered}
\]
Il codice usa quindi una costante e una soglia diverse da quelle
più precise indicate dopo \eqref{eq:kernel}. Se $x_j\ne0$,
$|y_j|\leq M$ e $\sum_jx_j^{-2}<\infty$, verifica la sommabilità
dei resti e la stima
\[
\left|\sum_j\left[R_{x_j,y_j}(\tau)
              +\frac{4\tau y_j}{x_j^2+\tau^2}\right]\right|
\leq\frac{36M^3}{\tau}\sum_jx_j^{-2}\longrightarrow0.
\]
È verificata anche la separazione delle serie convergenti che dà
\[
\sum_jR_{x_j,y_j}(\tau)
+4\tau\sum_j\frac{y_j}{x_j^2+\tau^2}\longrightarrow0.
\]
Questo controllo fornisce un maggiorante $O(\tau^{-1})$ della
somma dei resti, sufficiente al limite; non certifica il tasso
più forte $O(\tau^{-2})$ indicato nel testo.

Per la sostituzione delle parti reali è verificata la disuguaglianza
\[
\begin{gathered}
\left|\frac{\tau}{x^2+\tau^2}
       -\frac{\tau}{a^2+\tau^2}\right|
\leq D\left(\frac1{x^2}+\frac1{a^2}\right),\\
x,a,\tau>0,\quad |x-a|\leq D.
\end{gathered}
\]
Applicando formalmente la convergenza dominata alle serie, il
codice dimostra che la somma delle differenze, moltiplicate per
coefficienti di modulo al più uno, tende a zero quando gli
inversi dei quadrati di entrambe le successioni sono sommabili.
Per gli errori sulle altezze è formalizzato l'argomento che
separa una testa finita e una coda, assumendo pesi positivi e
sommabili, massa limitata e convergenza a zero di ogni singolo peso.

\paragraph{Obblighi non ancora verificati in Lean.}
Non sono ancora formalizzati l'enunciato geometrico completo
di \eqref{eq:target}, i risultati spettrali di Bianchi e la loro
applicazione ai rami sulla sfera, la fattorizzazione di Hadamard
con il passaggio al prodotto e ai logaritmi accoppiati, né i
limiti di Laplace ricavati dalle sezioni. Rimane soprattutto da
formalizzare il Lemma \ref{lem:roughcompetitor}: regolarità del
supporto, assenza di parte singolare della misura di area,
asintotica delle sezioni anche nei punti degeneri e recupero
delle identità sulle densità. Restano inoltre il collegamento
con i pesi specifici di \eqref{eq:universalmean}, la loro massa
limite e la conclusione mediante l'unicità di Minkowski.
I due lemmi centrali \ref{lem:roughcompetitor} e \ref{lem:mean}
non sono pertanto certificati nella loro interezza.

\paragraph{Ripetibilità e significato.}
Il comando \texttt{./check.sh}, eseguito dalla cartella
\texttt{formalizzazione}, ripete la verifica. Il relativo
\texttt{README.md} elenca ipotesi, corrispondenze e obblighi
residui; la sottocartella \texttt{verifica} conserva registri,
versioni e impronte dei sorgenti controllati.
Un lemma formalmente verificato rimane condizionato dalle ipotesi
nel suo enunciato: l'audit degli assiomi non dimostra tali ipotesi
per i corpi convessi. Il risultato qui registrato non autorizza
quindi a descrivere la prova candidata dell'unicità come
interamente verificata in Lean.

\begin{thebibliography}{9}
\bibitem{B16} G. Bianchi,
\emph{The covariogram and Fourier--Laplace transform in $\mathbb C^n$},
Proceedings of the London Mathematical Society 113 (2016), 1--23.
\href{https://arxiv.org/abs/1312.7816}{arXiv:1312.7816}.
\bibitem{AB07} G. Averkov, G. Bianchi,
\emph{Retrieving convex bodies from restricted covariogram functions},
Advances in Applied Probability 39 (2007), 613--629.
\href{https://arxiv.org/abs/math/0702892}{arXiv:math/0702892}.
\bibitem{B23} G. Bianchi, \emph{The covariogram problem},
in Harmonic Analysis and Convexity (2023), 37--82.
\href{https://people.dimai.unifi.it/bianchi/lavori/covariogram_survey14.pdf}{Manoscritto dell'autore}.
\bibitem{Siu24} Y.-T. Siu,
\emph{Weierstrass and Hadamard Factorization of Entire Functions},
note del corso Math 213a, Harvard University, autunno 2024.
\href{https://people.math.harvard.edu/~siu/math213a/weierstrass_and_hadamard_factorization_of_entire_function.pdf}{Note dell'autore}, p.~4.
\end{thebibliography}
\end{document}
